\subsubsection{Controlled Unitary Operators}\label{ex:controlled-unitary-operators}

In general, if we apply a unitary operator $U$ directly to one of its eigenstates:
\begin{equation*}
    U\ket{v} = e^{i\phi}\ket{v}
\end{equation*}
The eigenvalue only creates a \textbf{global phase}. Since global phases are physically unobservable, the phase information is effectively ``\emph{hidden}''. The natural question is therefore: \emph{\textbf{can we apply the unitary only to part of the quantum state, so that the phase becomes relative instead of global?}} The answer is yes, by using a \textbf{controlled unitary operator}.

\highspace
\begin{flushleft}
    \textcolor{Green3}{\faIcon{question-circle} \textbf{Motivation}}
\end{flushleft}
Suppose we have two qubits:
\begin{itemize}
    \item A \textbf{control qubit}, which decides whether the operation is executed;
    \item A \textbf{target qubit}, on which the unitary $U$ acts.
\end{itemize}
Instead of always applying $U$ to the target qubit:
\begin{equation*}
    \ket{v} \longrightarrow U\ket{v} = e^{i\phi}\ket{v}
\end{equation*}
We apply it \textbf{conditionally}:
\begin{itemize}
    \item If the control qubit is \ket{0}, do nothing to the target qubit;
    \item If the control qubit is \ket{1}, apply $U$ to the target qubit.
\end{itemize}
Graphically, this is represented as follows:
\begin{center}
    \begin{quantikz}
        \lstick{\ket{c}} & \ctrl{1} & \\
        \lstick{\ket{v}} & \gate{U} & 
    \end{quantikz}
\end{center}

\noindent
This simple modification is the key ingredient behind phase kickback.

\begin{definitionbox}[: Controlled Unitary]
    A \definition{Controlled Unitary Operator} is the two-qubit operator $CU$ that performs:
    \begin{equation*}
        \ket{0} \ket{\psi} \longrightarrow \ket{0} \ket{\psi}, \quad
        \ket{1} \ket{\psi} \longrightarrow \ket{1} U\ket{\psi}
    \end{equation*}
    Only the \textbf{target qubit} is \textbf{modified}. The \textbf{control qubit} itself is \textbf{never changed}.

    \highspace
    Since the operator acts on two qubits, its matrix has dimension $4 \times 4$. It has the block form:
    \begin{equation}\label{eq:controlled-unitary-matrix}
        CU = \begin{pmatrix}
            I & 0 \\
            0 & U
        \end{pmatrix}
    \end{equation}
    Where:
    \begin{itemize}
        \item The upper-left block corresponds to the control being \ket{0};
        \item The lower-right block corresponds to the control being \ket{1}.
    \end{itemize}
\end{definitionbox}

\highspace
\begin{flushleft}
    \textcolor{Green3}{\faIcon{book} \textbf{Action on Computational-Basis States}}
\end{flushleft}
Assume that the target is an eigenstate:
\begin{equation*}
    U\ket{v} = e^{i\phi}\ket{v}
\end{equation*}
The two possible values of the control qubit are:
\begin{itemize}
    \item \hl{Control qubit is $\ket{0}$}. Initially, the two-qubit state is:
    \begin{equation*}
        \ket{0}\ket{v}
    \end{equation*}
    Applying the controlled unitary, nothing happens:
    \begin{equation*}
        CU\ket{0}\ket{v} = \ket{0}\ket{v}
    \end{equation*}
    Graphically, this is represented as follows:
    \begin{center}
        \begin{quantikz}
            \lstick{\ket{0}} & \ctrl{1} & \\
            \lstick{\ket{v}} & \gate{U} & 
        \end{quantikz}
    \end{center}
    Since the control is \ket{0}, the ``\emph{switch}'' remains off.


    \item \hl{Control qubit is $\ket{1}$}. Initially, the two-qubit state is:
    \begin{equation*}
        \ket{1}\ket{v}
    \end{equation*}
    Applying the controlled unitary, the target is modified:
    \begin{equation*}
        CU\ket{1}\ket{v} = \ket{1} \otimes U\ket{v} = \ket{1} \otimes e^{i\phi}\ket{v} = e^{i\phi}\ket{1}\ket{v}
    \end{equation*}
    \textcolor{Red2}{\faIcon{exclamation-triangle}} Exactly the same eigenvalue $e^{i\phi}$ appears, but \textbf{nothing useful} is gained, since it is still a \textbf{global phase}.
\end{itemize}

\highspace
\begin{flushleft}
    \textcolor{Green3}{\faIcon{question-circle} \textbf{Why do we still need a Controlled Unitary?}}
\end{flushleft}
At this point, the controlled operator may seem disappointing because when the control is \ket{1}, the eigenvalue is still a global phase (irrelevant). However, notice something important: \hl{the phase appears \textbf{only when the control qubit is $\ket{1}$}.} This means that if the control qubit were somehow simultaneously in a superposition of \ket{0} and \ket{1}, the phase would affect only one branch of the computation.

\highspace
So, instead of choosing a simple computational-basis state for the control qubit, we prepare the control qubit in a superposition:
\begin{equation*}
    \ket{+} = \frac{\ket{0} + \ket{1}}{\sqrt{2}}
\end{equation*}
Now the controlled unitary acts only on the \ket{1} branch of the superposition:
\begin{equation*}
    CU\ket{+}\ket{v} = \frac{1}{\sqrt{2}} \left( CU\ket{0}\ket{v} + CU\ket{1}\ket{v} \right) = \frac{1}{\sqrt{2}} \left( \ket{0}\ket{v} + e^{i\phi}\ket{1}\ket{v} \right)
\end{equation*}
Now the phase no longer multiplies the \textbf{whole state}, but \textbf{only one branch} of the superposition. This transforms the previously invisible global phase into an observable relative phase.

\highspace
\textcolor{Green3}{\faIcon{check-circle} \textbf{Factor out $\ket{v}$.}} Both branches of the superposition contain the same $\ket{v}$, so we can factor it out:
\begin{equation*}
    CU\ket{+}\ket{v} = \frac{1}{\sqrt{2}} \left( \ket{0} + e^{i\phi}\ket{1} \right) \otimes \ket{v}
\end{equation*}
Now, the phase is \textbf{not} multiplying the target qubit anymore, instead $e^{i\phi}$ appears only in front of the \ket{1} component of the \textbf{control qubit}. This is the essential mathematical mechanism behind \textbf{phase kickback}. The next section will formalize this phenomenon and explain why we say the phase is ``\emph{kicked back}'' from the target qubit to the control qubit.